\section{Conclusion}
Neural Bellman Operators separate proposal, evaluation, improvement, and certification so that a computational policy can be used in an economic counterfactual. The substantive question is not whether a method contains a neural network, but whether its representation and its error account support the change in primitives being studied.

For rewards with fixed transitions, the policy-feature and optimistic-support geometry is established. We retain its complete finite-model implementation, action-sign exclusions, and repaired policies, while identifying precisely the part shared with scalar optimistic linear support. For changing transition laws, an endpoint chord alone is invalid. The cross-operator correction in Theorem \ref{thm:corrected_chord}, feasible Bernstein policy values, and statewise switching provide a continuum welfare guarantee that incorporates this missing effect. Approximate neural endpoints can enter through certified residual supersolutions rather than assumed exact evaluation.

In the stopped preference economy, operating benefits and liquidation levels reduce to a contract contrast. A risk-specific duration--effort ratio then characterizes regular movements of an indifference locus, and a matched restriction on deliberate preference adjustment identifies the option-value contribution to the ranking of risky positions. A primitive consumption-exposure condition supplies a positive risk-specific option bound over a probability--cost family. In the unchanged finite stopped economy, tensor Bernstein upper and lower values certify opposite risk rankings throughout a joint contract--transition region. This connects the reusable transition representation to the economic decision itself. The numerical threshold brackets retain their separate direct-solver interpretation; uniqueness and a diffusion interpretation are not inferred from a strict regional sign certificate.

The matched count comparison determines what the transition correction buys. It is not the sharpest upper oracle: localization with the same policy bank favors count information. Its gain is compressed horizon work and storage, with a total coefficient-error bound that allows policy kinks. The common-bank frontier and the autonomous refinement experiment separately report the price of precision and the demand for new endpoint solves. Lower-policy storage and exhaustive finite-action work remain in that account.

The resource comparison establishes the importance of reusable structure on both sides of a computational comparison. A nonlinear convex critic improves on the fitted quadratic continuation in the recorded experiments, but an exact reusable parametric quadratic program is faster than learned deployment for the tested short quadratic economy. Neither that fact nor the mesh-only policy-bank ablation invalidates the transfer certificate. They determine what can be attributed to learning and what should instead be attributed to policy evaluation, structural optimization, and parameter reuse.

The new economic analysis makes the decision horizon and cross-preference cardinal primitive explicit. A first-risk choice concerns a positive-duration commitment, not a control changed at a single instant. The primitive option result survives a nontrivial common-tilt family, while full-model counterfactuals show that its dynamic magnitude and its effect on risk choice depend on preference-state payoffs, liquidation, duration, and covariance. Wealth-lottery perturbations move the original indifference boundary; their failures remain alongside an additional common-region certificate. These are stronger distinctions between economic propositions, not substitutes for the target-error control needed to transfer a sign to the diffusion benchmark.

At the original finite decision target, streaming chord and count implementations certify the same risk comparison with different work. The experiment charges common construction and validation costs, retains a mesh-only lower-policy comparator, and reports kernel, policy, and coefficient storage separately. It therefore links compression to the economic accuracy requirement without converting a component saving into a claim about neural necessity or total process memory.

The continuous-time formulation, finite transition targets, neural proposals, recursive preferences, temporal selves, and persistent dynamic competition remain distinct parts of the framework. Keeping their assumptions and evidence explicit permits new representations and counterfactuals without changing the economic meaning of the reported welfare guarantee.
